\chapter{并行与分布式计算}

\section{领域导言与演进背景}

\begin{preamblebox}[知识领域概述与变革]
并行与分布式计算 (PDC)致力于组织、协调与控制跨越不同物理时空且并发发生的多种计算活动。伴随处理器内部晶体管规模的激增、单机中多核处理器的普及以及网络化计算节点的无处不在，并行性与分布式特性已成为现代计算系统的本质属性。相较于传统串行程序，并行与分布式架构不仅能够成倍释放硬件潜能、大幅提升计算性能，更是应对现实物理世界固有互联性、部件局部失效以及通信不可靠性的必然技术路径。通过打破传统串行编程中“所有计算步骤必须在单一位置按严格时序依次发生”的桎梏，并行与分布式计算揭示出贯穿现代计算系统各个层级的全套核心抽象、设计范式与分析方法。

在学术界与工业界的通用表述中：
\begin{itemize}
    \item \textbf{并行编程}：侧重于调度和组织多个计算活动在时间上的“协同并发发生”；
    \item \textbf{分布式编程}：侧重于协调发生在不同物理“空间位置”上的独立计算活动；
    \item \textbf{并发编程}：侧重于管理正在推进中的多项活动之间及其与外部环境的“交错交互”。
\end{itemize}
然而在现实工程中，这三者往往深度交织、密不可分。并行性天然蕴含某种程度的分布性，因为除依赖时间片上下文切换或量子效应外，无顺序依赖的多个并发操作必然发生于不同的物理硬件单元之上；反之，在缺乏明确通信约束时，发生于不同地理位置的操作彼此之间亦不存在绝对的先后次序。本知识领域既重点聚焦“主动显式引入并行与分布式”的技术语境，也涵盖“因数据或服务天然跨地域分布而不得不引入并发控制”的被动约束场景。

将并行与分布式计算讲深讲透极具教学挑战性。该领域已由早期的前沿高级选修课集群，演进为渗透至现代计算学科几乎每一个角落的通用核心知识体系。现实中，几乎所有具备串行解法的问题都存在并行或分布式解法，而更多前沿问题唯有在并发语境下才会真正涌现。历史上，该领域的各分支在不同计算子领域中呈现出非均衡的碎片化演进，有时甚至使用截然不同的术语、工程实践和课程载体，从而遮蔽了底层基本思想与核心技能的高度一致性。本知识领域的核心使命，正是提炼并讲授这一横跨软硬件体系的坚实知识基底。

\textbf{自 CS2013 以来的核心变革：}
\begin{enumerate}
    \item \textbf{知识体系抽象重构与原语凝练}：摒弃以往简单罗列特定多线程库（如 Pthreads、OpenMP）、消息传递库（如 MPI）或特定集群中间件（如 Hadoop）的做法，重塑为由“程序与执行”、“通信机制”、“协调与同步”、“软件工程”、“算法与应用领域”五大核心抽象原语构成的现代知识图谱。
    \item \textbf{精炼解耦的计算机核心}：将计算机专业通识必修的核心学时精炼为 10 学时（每个知识单元各 2 学时），聚焦于破除单机串行顺序执行假设的基本规律、Happens-before 偏序模型及非确定性并发直觉，不捆绑任何特定语言或特定硬件架构。
    \item \textbf{弹性模块化的领域核心}：领域核心（共 26 学时）采用“选修以下一项或多项”的拓展机制，允许各高校根据自身师资与专业定位，灵活构建高性能计算（HPC）、异构加速（GPU/SIMD）、多核共享内存系统、弹性云原生或去中心化容错存储等多元课程方案。
    \item \textbf{全面引入真实能力规范}：设置真实工程任务，重点考察学生在多核/异构/分布式环境下的并发重构、性能剖析（Work-Span 跨度分析与 Amdahl 定律应用）、死锁预防以及容错共识协议设计能力。
\end{enumerate}

\textbf{本领域课程起草委员会：}
\begin{itemize}
    \item \textbf{主席}：Doug Lea（纽约州立大学奥斯威戈分校）
    \item \textbf{委员}：Sherif Aly（开罗美国大学）、Michael Oudshoorn（高点大学）、项乔（Qiao Xiang，厦门大学）、Dan Grossman（华盛顿大学）、Sebastian Burckhardt（微软研究院）、Vivek Sarkar（佐治亚理工学院）、Maurice Herlihy（布朗大学）、Sheikh Ghafoor（田纳西理工大学）、Chip Weems（马萨诸塞大学阿默斯特分校）
    \item \textbf{主要贡献者}：Paul McKenney、Peter Buhr
\end{itemize}
\end{preamblebox}

\section{核心学时分配建议}

本知识领域在“向日葵模型”下的学时推荐如下表所示。计算机核心涵盖所有计算专业学生打破串行思维所必备的通识基础，领域核心则提供进阶深度选修指南：

\begin{table*}[htbp]
\centering\small
\setlength{\tabcolsep}{2.5pt}\footnotesize\begin{tabularx}{0.88\textwidth}{Y c c}
\toprule
\textbf{知识单元} & \textbf{计算机核心} & \textbf{领域核心} \\
\midrule
程序与执行 & 2 & 2 \\
通信机制 & 2 & 6 \\
协调与同步 & 2 & 6 \\
软件工程 & 2 & 3 \\
算法与应用领域 & 2 & 9 \\
\midrule
\textbf{合计} & \textbf{10 学时} & \textbf{26 学时} \\
\bottomrule
\end{tabularx}
\end{table*}

\textbf{学时说明与实施建议：}
\begin{itemize}
    \item \textbf{计算机核心}：共 10 学时，覆盖并行与分布式计算的核心理念。其主要教学目标是消除严格串行编程的思想禁锢，揭示并行与分布式系统及软件的本质结构与核心特性。学习成果侧重于开发小规模多活动并发程序，并分析其基础性质。该学时指标不包含特定编程语言、开发工具、计算框架与物理平台的具体细节教学（这些内容通常由后续具体专业课程结合实验项目承载）。
    \item \textbf{领域核心}：共 26 学时，各单元内容采用“选修以下一项或多项”的形式组织，为不同侧重点的专业选修课提供弹性空间。例如：
    \begin{itemize}
        \item 高性能计算（HPC）方向：聚焦异构数据并行与线性代数求解；
        \item 计算机系统方向：聚焦多核共享内存协调、中间件与系统资源管理；
        \item 分布式大数据方向：聚焦容错存储、分布式共识与大规模云服务。
    \end{itemize}
\end{itemize}

\section{知识单元与详细大纲}

\subsection{程序与执行}
\textbf{【学时建议】}：2 计算机核心学时，2 领域核心学时。

本单元奠定并行与分布式计算的核心抽象，重点讲授并行动作的定义、执行拓扑、时序关系以及启动控制机制。

\subsubsection{计算机核心主题}
\begin{itemize}
    \item \textbf{基本定义与核心属性}：
    \begin{itemize}
        \item \textbf{可并行化动作}：包括闭包、函数调用以及网络服务调用；复合动作、会话、并发任务与作用域；动作的实体命名与标识（如线程 ID、进程号、URI）；计算粒度与调度开销权衡。
        \item \textbf{动作先后定序与先行发生关系}：放宽严格串行时序约束；用串行/并行有向无环图建模程序组件的非确定性执行流动。
        \item \textbf{独立性判定}：基于操作的可交换性、数据流依赖性与前置条件，严格判定何时执行顺序互不相干。
        \item \textbf{一致性约定}：关于数值状态与谓词的协同协议；竞态条件、原子性与分布式共识直觉。
        \item \textbf{并行程序中的必要定序保障}：互斥锁、内存屏障、安全发布；实体间通信施加的时序制约（例如：消息的“发送”必然发生在“接收”之前）。
        \item \textbf{执行位置与执行实体}：承载并行动作的物理设备（硬件核、远程主机）；单一执行位置对多个并发实体的复用与调度；外部世界及人类用户作为并发实体的交互模式；基于时间片分时调度与虚拟化机制在少量处理器上仿真大量并行动作。
        \item \textbf{实体失效与通信故障}：不可信实体及不受本程序或管理域控制的网络协议所引发的系统失效；容错级别作为系统核心设计维度的折中考量。
    \end{itemize}
    \item \textbf{并行动作的启动机制}：
    \begin{itemize}
        \item \textbf{编排放置}：协调动作最终由指定物理实体或节点执行（从硬编码静态绑定、配置脚本到现代云平台的自动化弹性编排）；建立端到端通信与资源分配。
        \item \textbf{过程式启动}：在特定程序断点处派生并发动作（如创建新线程、Fork 子进程）；利用结构化作用域确保所有衍生动作全部结束时作用域方能退出。
        \item \textbf{反应式/事件驱动启动}：由外部事件触发并行动作（如注册事件回调处理器）；对动作具体启动和终止时间的非精确控制。
        \item \textbf{依赖驱动启动}：基于前序计算或前置任务的全部完结状态触发后续任务（如流式计算、延续触发 Continuations）。
    \end{itemize}
\end{itemize}

\subsubsection{领域核心与拓展主题}
\begin{itemize}
    \item \textbf{底层硬件映射与物理执行机制（选择以下一项或多项）：}
    \begin{itemize}
        \item CPU 数据级与指令级并行（ILP、流水线超标量、SIMD 向量扩展）。
        \item SIMD 向量化架构与异构数据并行（GPU、专用深度学习加速芯片）。
        \item 多核处理器中的抢占/协作调度并发、任务池与 Actor 消息传递模型。
        \item 计算集群、云计算基础设施与弹性算力资源编排。
        \item 参与者无边界、跨管理域的开放分布式系统。
        \item 新兴前沿计算模型：量子并行计算、DNA/分子计算。
    \end{itemize}
\end{itemize}

\subsubsection{示范性学习成果与能力目标}
\begin{enumerate}
    \item \textbf{[应用 (计算机核心)]} 针对复杂的复合数值表达式（例如 $a = (b + c) \times (d + e)$），绘制其并行计算依赖有向无环图（DAG），展示操作间的并发度与关键路径。
    \item \textbf{[评估 (计算机核心)]} 在给定的并发代码片段中准确排查并指出数据竞态隐患。
    \item \textbf{[解释 (计算机核心)]} 结合具体工程场景，深入阐述将既有串行程序并行化在多大程度上能够提升系统吞吐量或降低端到端时延，并客观评估由此对系统整体能效与热功耗产生的实际影响。
    \item \textbf{[开发 (领域核心)]} 编写一个高效统筹高并发事件（如传感器密集采样或网络数据包接收）的高性能计数统计函数。
    \item \textbf{[开发 (领域核心)]} 分别采用多种不同的编程风格（如基于线程池、函数式流式处理、Actor 消息传递等）实现经典 Filter-Map-Reduce（过滤-映射-规约）程序。
    \item \textbf{[开发 (领域核心)]} 编写一个并发网络服务程序，通过为每个新接入的客户端派生独立线程（或其他过程式激活机制）处理并返回所请求的 Web 页面。
\end{enumerate}

\subsection{通信机制}
\textbf{【学时建议】}：2 计算机核心学时，6 领域核心学时。

本单元系统剖析并发与分布式实体间交换信息的基础介质、信道抽象、内存一致性以及数据存储模型。

\subsubsection{计算机核心主题}
\begin{itemize}
    \item \textbf{通信基础要素}：
    \begin{itemize}
        \item \textbf{通信介质与形态}：管道/通道（Channels，消息传递或 I/O 流）、共享内存、异构显存、结构化数据存储；对底层物理硬件、网络拓扑、互联协议、语言级运行时以及仿真模拟机制的依赖。
        \item \textbf{通信通道}：显式端到端实体通信；寻址与命名机制；核心 API 体系（网络套接字 Sockets、基于体系结构或语言原生的通道构造如 Go Channel）。
        \item \textbf{内存通信模型}：实体仅通过指定物理/虚拟地址直接与内存进行信息交换的架构；内存一致性基础（按位原子性边界、缓存一致性 Cache Coherence、局部内存读写定序）；内存层次结构与空间/时间局部性（高速缓存 Cache、访问时延、伪共享 False Sharing 问题）；跨异构存储体的数据搬运（如 CPU 主存与 GPU 显存间的显式 DMA 传输）；多层级数据共享域、作用域与私有缓存。
        \item \textbf{协同数据存储}：由多个并发实体协同维护的结构化数据（映射表 Maps、集合 Sets 等并发抽象数据类型）；存储模式形态（独占式 Owned、共享式 Shared、分片式 Sharded、副本复制式 Replicated、不可变数据 Immutable、多版本并发控制 Versioned）。
    \end{itemize}
    \item \textbf{基于通信的并发编程范式}：
    \begin{itemize}
        \item 熟练运用网络通道、套接字（Socket）以及远程过程调用（RPC）API 构建分布式通信组件。
        \item 掌握特定编程语言中基于共享内存模型的安全数据交换构件与内存屏障用法。
    \end{itemize}
\end{itemize}

\subsubsection{领域核心与拓展主题}
\begin{itemize}
    \item \textbf{高级通信属性与扩展机制（选择以下一项或多项）：}
    \begin{itemize}
        \item \textbf{通信介质深度特性}：
        \begin{itemize}
            \item 网络拓扑结构：单播、组播、信箱、交叉开关矩阵与自适应路由算法。
            \item 并发交互属性：消息定序、传递一致性保证、操作幂等性、计算与通信重叠技术。
            \item 可靠性工程：物理信道噪声、传输校验错误、丢包重传与滑动窗口机制。
            \item 数据封送与序列化：跨平台数据表示格式（Protobuf、JSON）、数据封送。
            \item 协议设计理论：活性保障、通信死锁规避。
            \item 安全性保障：传输保密性、数据完整性、端到端身份认证与授权机制。
            \item 性能评估维度：端到端时延、有效带宽与吞吐量、信道资源争用、系统响应活性。
            \item 排队论在信道拥塞与性能容量预测中的建模应用。
        \end{itemize}
        \item \textbf{通道高级控制与策略}：
        \begin{itemize}
            \item 通道管理策略：通信端点生命周期、会话保持、队列缓冲容量、饱和应对策略（阻塞等待 vs 丢弃丢包）、流量与速率控制。
            \item 发送控制（通常为过程式）与接收控制（通常为响应式回调或 RPC 服务派发）的编程对称性分析。
            \item 数据打包、序列化反序列化、输入合法性校验、传输加密与数据压缩。
            \item 海量慢速 I/O 设备与高并发连接下的多路复用与解复用机制：基于完成态驱动与调度器驱动的异步技术；async/await 异步流、select 与 poll/epoll 事件轮询 API。
            \item 进程通信形式化理论与代数分析（例如通信顺序进程 CSP 模型）。
        \end{itemize}
        \item \textbf{高级共享内存体系}：
        \begin{itemize}
            \item 现代语言与硬件内存模型：顺序一致性 (SC)与获取-释放一致性。
            \item 高并发内存管理：跨核心共享数据的安全内存回收机制；引用计数与替代方案（如 Hazard Pointers、基于历元 Epoch 的回收）。
            \item 块级数据布局与批量迁移：减少跨节点通信流量与提升局部性；计算与数据搬运的异步流水化重叠；数据内存布局对缓存命中的颠覆性影响（结构体数组 AoS 与 数组结构体 SoA 的性能对比）。
            \item 共享内存仿真：分布式共享内存（DSM）、远程直接内存访问（RDMA）。
        \end{itemize}
        \item \textbf{分布式数据存储与共识}：
        \begin{itemize}
            \item 一致性模型谱系：原子性、线性一致性、事务可串行化、缓存一致性、因果一致性、并发冲突消解、最终一致性以及区块链账本技术。
            \item 节点故障与网络分区：多数派法定人数表决；经典容错分布式共识协议（Paxos、Raft）。
            \item 安全性与零信任：拜占庭容错机制（BFT）、工作量证明（PoW）及新型权益证明共识协议。
        \end{itemize}
    \end{itemize}
\end{itemize}

\subsubsection{示范性学习成果与能力目标}
\begin{enumerate}
    \item \textbf{[评估 (计算机核心)]} 针对给定的工业应用场景与系统架构约束，全面论证并判定采用共享内存模型还是消息传递模型更具性能与工程实现优势。
    \item \textbf{[应用 (计算机核心)]} 编写一个经典的生产者-消费者程序（其中一个组件生成数据，另一个组件计算其移动平均值）；分别在小规模标量数据和大尺寸高精度数值场景下测量系统的实际加速比。
    \item \textbf{[开发 (领域核心)]} 编写一个能够将大规模数据集的不同分片并发派发给多个 Worker 计算节点并汇聚各局部运算结果的分布式程序（例如分布式超大数组求和）。
    \item \textbf{[开发 (领域核心)]} 编写一个并发数据汇聚程序，从多个异地远程站点并发拉取异构数据并利用某种形式的分布式规约输出汇总统计。
    \item \textbf{[评估 (领域核心)]} 编写基准测试，精确对比有缓冲队列与无缓冲同步通道在不同并发压力下的吞吐量表现与延迟瓶颈。
    \item \textbf{[评估 (领域核心)]} 针对给定的业务通信方案，全面评估其在防窃听、抗重放攻击与身份校验等维度是否具备充足的安全保障。
    \item \textbf{[分析 (领域核心)]} 构造一个包含数据竞态的并发内存访问代码序列，并证明该执行历史不满足顺序一致性。
    \item \textbf{[分析 (领域核心)]} 针对网络消息通信场景，构造一个由同步阻塞式发送操作相互等待而导致系统死锁的经典对抗案例。
    \item \textbf{[应用 (领域核心)]} 阐述至少一种在涉及多个互斥锁获取时杜绝死锁与活性故障的通用设计模式（如固定全局加锁顺序、两阶段锁）。
    \item \textbf{[开发 (领域核心)]} 编写一段实验程序，利用多线程交错执行直观展示硬件层面的内存乱序重排或消息异步乱序到达现象。
    \item \textbf{[评估 (领域核心)]} 深入剖析在不同数据更新写争用频率下，乐观并发控制（OCC）与悲观并发控制（PCC）的相对优劣与吞吐量临界点。
    \item \textbf{[分析 (领域核心)]} 举例构造一个由于高并发写冲突严重导致乐观更新尝试无限重试、甚至永远无法提交成功（活锁或饥饿）的极端场景。
    \item \textbf{[开发 (领域核心)]} 重构既有的并发软件系统，使其平滑迁移至具备更高可扩展性、更高容错性与更高可用性的现代分布式数据存储引擎之上。
\end{enumerate}

\subsection{协调与同步}
\textbf{【学时建议】}：2 计算机核心学时，6 领域核心学时。

本单元聚焦多个并发实体在执行时序上的相互制约与协同推进，涵盖任务依赖、推进活性、原子性互斥以及高级同步机制。

\subsubsection{计算机核心主题}
\begin{itemize}
    \item \textbf{协调与同步基础}：
    \begin{itemize}
        \item \textbf{依赖性动作的执行控制}：当某个活动的启动或推进依赖于其他活动时的执行调度与控制；基于完成态的控制（屏障 Barriers、线程汇合 Joins）；基于数据就绪驱动的控制（生产者-消费者设计模式）；基于条件满足的控制（忙等待轮询 Polling、重试 Retrying、指数退避 Backoffs、协同协助 Helping、挂起等待 Suspension、等待队列排队、条件信号唤醒 Signaling、超时时限控制 Timeouts）；响应式控制（激活并触发 Continuation 延续动作）。
        \item \textbf{推进活性与无死锁保障}：循环等待依赖与死锁；系统推进条件的单调性。
        \item \textbf{原子性语义与互斥锁}：底层硬件原子指令（CAS、LL/SC、Fetch-and-Add）；强制局部内存访问次序；互斥锁与临界区互斥；死锁防范工程技术（锁层次全局定序、锁粒度粗化、随机退避重试、利用锁管理器 Lock Manager 进行安全封装）；并发常见严重缺陷（遗漏加锁或未配对解锁、持有锁时调用不受控的未知外部操作、嵌套加锁死锁）。
    \end{itemize}
    \item \textbf{协调与同步编程实践}：
    \begin{itemize}
        \item 掌握并发组件优雅退出与协同终止控制机制。
        \item 熟练运用特定语言中的互斥锁、屏障、信号量等经典同步原语，在彻底杜绝数据竞态的同时保障系统的活性。
        \item 掌握在特定软件框架中利用事务性 API（如软件事务内存 STM 或分布式事务接口）实现高可靠并发协作。
    \end{itemize}
\end{itemize}

\subsubsection{领域核心与拓展主题}
\begin{itemize}
    \item \textbf{高级协调机制与控制流交互（选择以下一项或多项）：}
    \begin{itemize}
        \item \textbf{推进活性进阶理论}：非阻塞算法活性等级（无锁 Lock-free、无等待 Wait-free、无障碍 Obstruction-free）；调度公平性与优先级调度反转规避；同步机制与数据一致性、系统可靠性之间的权衡；性能瓶颈剖析（锁争用开销、锁粒度、锁护航 Convoying 效应、可扩展性极限）；典型无锁并发数据结构与算法实现（如 Treiber 栈、Michael-Scott 队列）。
        \item \textbf{原子性与资源控制进阶}：数据所有权与无锁借用语义；高级锁变体（顺序锁 SeqLocks、读写锁 Read-Write Locks、可重入锁、票据锁 Ticket Locks）；基于事务的并发控制（乐观与悲观事务执行）；分布式锁机制与租约可靠性管理。
        \item \textbf{与其他程序控制流形态的交互}：
        \begin{itemize}
            \item 屏障替代方案：物理时钟与逻辑时钟（Lamport 逻辑时钟、向量时钟 Vector Clocks、虚拟时间）；数据流模型与延续传递风格（CPS）；期约机制与异步远程调用；基于共识驱动的状态机复制；规约器与收集器结果汇聚机制。
            \item 投机执行、分支裁决与协同取消：投机执行对微架构安全与系统可观测性的潜在风险（如侧信道漏洞）。
            \item 基于系统资源的同步：经典计数信号量与条件变量。
            \item 控制流拓扑：细粒度任务调度器、带动态选主的串行-并行混合循环、流水线与流式处理管道、嵌套分治并行。
            \item 异常与故障协同响应：分布式异常传播与捕获处理、心跳监测超时探测、自愈容错与多数派故障仲裁机制。
        \end{itemize}
    \end{itemize}
\end{itemize}

\subsubsection{示范性学习成果与能力目标}
\begin{enumerate}
    \item \textbf{[应用 (计算机核心)]} 针对一段包含典型同步缺陷或数据竞态的并发代码，展示如何引入合理的互斥机制将其修复为线程安全程序。
    \item \textbf{[开发 (计算机核心)]} 编写一个协调程序，确保当且仅当一组并发异步子任务全部安全完成计算后，主控程序方能正常关闭并优雅终止。
    \item \textbf{[开发 (领域核心)]} 实现一个能够高吞吐、低延迟安全聚合海量并发传感器输入或网络数据包计数的高性能非阻塞计数器。
    \item \textbf{[开发 (领域核心)]} 分别运用不同的并发协调范式（如屏障汇聚法、Actor 消息归约法）完整实现 Filter-Map-Reduce 处理引擎。
    \item \textbf{[开发 (领域核心)]} 编写一段多阶段流水化并行程序，严格保障前一组并行任务全部汇聚提交之后，后一组并行任务才能被触发分发。
    \item \textbf{[开发 (领域核心)]} 编写一个采用投机执行机制的并发解题搜索程序，一旦其中某个工作线程率先检索到有效答案，能够立即广播信号并安全取消终止其他所有线程的无谓搜索。
    \item \textbf{[开发 (领域核心)]} 编写一段健壮的并行协同处理程序，当某个并发执行流突发致命数值异常（如除以零或浮点溢出）时，能向所有协作实体广播故障信号并协同安全回滚退出。
    \item \textbf{[评估 (领域核心)]} 编写一个使多个分布式对等节点就当前统一时钟达成共识的简易协议，并深入剖析其与工业级标准网络时间协议（NTP）相比在网络抖动和时钟漂移下的本质局限。
    \item \textbf{[开发 (领域核心)]} 实现一个轻量级高并发网络服务器，为每个接入的新客户端创建独立并发上下文（如轻量级协程或虚拟线程）处理并响应 Web 请求。
\end{enumerate}

\subsection{并发与分布式软件工程}
\textbf{【学时建议】}：2 计算机核心学时，3 领域核心学时。

本单元系统讲授并发与分布式软件在规约制定、质量属性保障、缺陷排查以及性能评估方面的软件工程方法论。

\subsubsection{计算机核心主题}
\begin{itemize}
    \item \textbf{安全性、活性与性能指标需求}：
    \begin{itemize}
        \item 运用时态逻辑基础概念严格界定系统的核心不变量：区分“坏事绝不会发生”与“好事终将发生”。
        \item 关键工程度量指标体系：吞吐量、响应性、端到端延迟、服务可用性、系统能耗与碳足迹、水平/垂直可扩展性、硬件资源利用率、通信带宽开销、排队等待时间与速率限制、调度公平性指标、工业服务等级协议（SLA）指标规范。
    \end{itemize}
    \item \textbf{违规识别、测试验证与缺陷修复}：
    \begin{itemize}
        \item 并发软件常见故障分类识别：时序定序缺失（竞态错误 Race Conditions）、原子性破裂（如典型的检查-执行竞争 Check-Then-Act 漏洞）、非终止循环（活锁 Livelock、死锁 Deadlock）。
    \end{itemize}
\end{itemize}

\subsubsection{领域核心与拓展主题}
\begin{itemize}
    \item \textbf{并发与分布式软件的规范化与评测（选择以下一项或多项）：}
    \begin{itemize}
        \item \textbf{形式化规约与建模}：串行规约向并发语义的拓展（如线性一致性 Linearizability 规范）；协议交互契约、会话类型与跨节点事务规范；借助工业级形式化工具（如 UML 交互图、TLA+ 规范语言、程序逻辑推演）验证并发协议的不变量；恶意节点与拜占庭行为下的安全性与活性证明；通信通道隔离性要求（防跨层侧信道信息泄漏）、入口流量过滤与反 DoS 速率限制。
        \item \textbf{静态分析与算法建模}：面向并发正确性、吞吐量瓶颈、时延抖动、资源争用与能耗分析的静态代码审查技术；基于 DAG 图模型的并行算法效能分析（总工作量 Work $T_1$、关键路径跨度 Span $T_\infty$、理论并行度 $T_1/T_\infty$ 与贪心调度边界）。
        \item \textbf{经验性评估与动态测试}：并发测试与调试专用工具链（数据竞态动态检测器如 ThreadSanitizer、模糊测试工具 Fuzzers、锁依赖环检测工具、高负载压力测试与严苛环境折磨测试 Torture Tests、集成至 CI/CD 流水线的并发自动化测试生成器）；经验性多维度量化评测（实测并深度对比吞吐瓶颈、同步额外开销、线程阻塞等待、通信流量开销、数据局部性破坏以及突发流量/海量并发连接冲击下的系统退化表现）；针对特定应用领域的专有评测分析与调优技术。
    \end{itemize}
\end{itemize}

\subsubsection{示范性学习成果与能力目标}
\begin{enumerate}
    \item \textbf{[开发 (计算机核心)]} 在不破坏既有系统核心业务语义、不变量与核心功能的前提下，重构现有软件规约以引入并行处理与分布式解耦。
    \item \textbf{[解释 (计算机核心)]} 深入阐述为何在并发与分布式软件架构设计中，严格规避“非局部副作用”是一条不可动摇的黄金设计法则。
    \item \textbf{[应用 (计算机核心)]} 针对一个批量并行计算过程，精确指明并形式化定义在每个并发计算步骤中必须严格维持的一组状态不变量。
    \item \textbf{[开发 (计算机核心)]} 编写一个高灵敏度的自动化测试程序，能够高概率稳定暴露潜在的数据竞态错误（例如两个并发活动在未加同步的情况下同时对共享累加器执行递增操作导致写丢失）。
    \item \textbf{[评估 (领域核心)]} 针对某在线服务突发超出设计承载能力的超大规模请求洪峰，精确制定系统优雅降级规约并实测其过载防护表现。
    \item \textbf{[开发 (领域核心)]} 准确诊断并重构消除某分布式应用中由串行执行瓶颈（Amdahl 定律瓶颈）导致的性能瓶颈，使其重获线性加速能力。
    \item \textbf{[评估 (领域核心)]} 借助现有的工业级测试基准脚手架，实测并全面对比针对同一接口设计的两种不同并发底层实现版本在多核高并发下的吞吐效能与锁竞争损耗。
    \item \textbf{[开发 (领域核心)]} 诊断并有效修复因跨节点网络通信过于频繁或数据读取延迟过高而导致的系统整体吞吐暴跌故障（如引入数据预取、批量聚合与异步流水化）。
    \item \textbf{[开发 (领域核心)]} 准确定位并优化并发系统中因过度的线程频繁创建销毁、锁争用严重或小对象内存分配而导致的系统性能抖动与资源开销膨胀。
    \item \textbf{[开发 (领域核心)]} 针对分布式系统中单点宕机、网络偶发丢包或超时等典型故障，精准诊断并修复系统的可靠性或高可用性脆弱点。
\end{enumerate}

\subsection{并行算法与应用领域}
\textbf{【学时建议】}：2 计算机核心学时，9 领域核心学时。

本单元系统梳理主流并行与分布式计算范式，并深入探讨经典并行算法在各个核心工业与科研领域的落地实践。

\subsubsection{计算机核心主题}
\begin{itemize}
    \item \textbf{并行与分布式算法的表达与实现}：熟练掌握在现代主流编程语言与开发框架中落地并行与分布式算法的核心技能，包括派生并发执行活动（如线程、纤程、Actor）、安全访问共享内存数据结构，以及熟练调用通道、套接字（Socket）与远程过程调用（RPC）网络接口。
    \item \textbf{核心计算范式与系统架构全景速查}：
\end{itemize}

\begin{table*}[htbp]
\centering\small
\begin{tabularx}{\textwidth}{p{0.14\textwidth} p{0.16\textwidth} Y Y p{0.18\textwidth}}
\toprule
\textbf{系统范式分类} & \textbf{典型执行实体} & \textbf{典型通信机制} & \textbf{典型算法领域} & \textbf{核心工程目标} \\
\midrule
\textbf{多核系统} & 操作系统线程、轻量级任务 & 共享内存、硬件原子指令、互斥锁 & 系统资源调度、高并发实时数据处理 & 极致吞吐量、低时延、高能效比 \\
\midrule
\textbf{反应式系统} & 事件处理器、异步回调函数 & 非阻塞 I/O 通道、事件总线 & 在线微服务网关、实时交互控制系统 & 超低响应时延、高可用性 \\
\midrule
\textbf{数据并行系统} & GPU 核心、SIMD 向量单元、AI 加速器 & 异构内存空间、显存 DMA 批量搬运 & 稠密线性代数、计算机图形渲染、大数据分析 & 极致算力吞吐、能耗效率 \\
\midrule
\textbf{计算集群} & 托管物理计算节点、MPI 进程 & 高速以太网/InfiniBand 套接字、消息通道 & 复杂科学数值仿真、大规模分布式图挖掘 & 大规模批处理吞吐量 \\
\midrule
\textbf{云计算平台} & 弹性容器、虚拟机、Serverless 函数 & RESTful/gRPC 服务接口、分布式消息总线 & 大规模 Web 互联网应用、海量用户在线服务 & 弹性伸缩性、故障自动隔离 \\
\midrule
\textbf{开放分布式系统} & 跨地域跨管理域自治节点、P2P 节点 & 网络套接字、分布式去中心化数据存储 & 拜占庭容错存储、区块链共识系统 & 强容错性、去中心化信任 \\
\bottomrule
\end{tabularx}
\end{table*}

\subsubsection{领域核心与拓展主题}
\begin{itemize}
    \item \textbf{八大核心算法应用领域（选择以下一项或多项）：}
    \begin{itemize}
        \item \textbf{科学计算与线性代数}：向量与高维矩阵运算（BLAS 算子级并行）；数值计算精度与数值稳定性控制；线性代数在现代大数据分析、深度神经网络大模型训练与推理中的并行加速应用。
        \item \textbf{数据处理与检索}：并行排序（如双调排序 Bitonic Sort、并行样本排序）、分布式哈希查找与倒排索引检索、高吞吐并发数据结构设计。
        \item \textbf{图算法、搜索与组合优化}：并行图遍历与节点标记、大规模边并行化处理、搜索分支剪枝与限界分析、投机搜索并发化、基于海量真实网络拓扑的分布式图计算（如 PageRank、连通分量求解）。
        \item \textbf{建模与物理系统仿真}：常微分与偏微分方程数值并行解法（有限元法、有限差分法）；蒙特卡洛随机模拟；高并发 N 体问题天体引力仿真；遗传算法与群体智能演化计算。
        \item \textbf{形式逻辑与约束求解}：大规模并行命题逻辑可满足性求解器、并发逻辑程序设计范式。
        \item \textbf{计算机图形学与计算几何}：大规模坐标仿射几何变换、图形光栅化渲染管线、实时硬件光线追踪加速、高精度网格碰撞检测。
        \item \textbf{算力与系统资源管理}：CPU/GPU 处理单元、内存带宽、通信通道与主机节点的动态分配、空间放置、内存回收与负载均衡调度；独占资源与共享资源的并发调度冲突消解；静态排程、动态调度与弹性伸缩算法；任务批量合并、优先级反转规避、资源隔离分区；基于工作窃取去中心化任务队列的高性能负载均衡机制。
        \item \textbf{网络与金融服务系统}：高并发分布式 Web API 网关架构、电子加密货币协议、高频交易引擎与分布式 ACID 事务处理系统、大规模多人在线交互系统（如 MMO 游戏）的分布式状态同步引擎。
    \end{itemize}
\end{itemize}

\subsubsection{示范性学习成果与能力目标}
\begin{enumerate}
    \item \textbf{[开发 (计算机核心)]} 基于已知的经典并行算法（如并行归并、分治前缀和求值、并行 BFS），使用现代并发库开发出高可靠的并发/分布式组件。
    \item \textbf{[开发 (计算机核心)]} 编写一个数据并行程序（例如利用 SIMD 指令或 OpenMP 指令集并行计算百万级浮点数数组的均值与方差）。
    \item \textbf{[开发 (计算机核心)]} 对既有的单线程事件驱动串行程序进行多线程化改造，在事件回调函数（如 GUI 界面交互事件或网络请求处理器）中派生异步工作线程，彻底消除主事件循环的卡顿与阻塞。
    \item \textbf{[开发 (计算机核心)]} 通过识别计算瓶颈并引入并行化或分布式扩展技术，显著提升某个性能受限串行软件模块的实际执行效能。
    \item \textbf{[评估 (计算机核心)]} 针对给定的实际系统工程诉求，在多核共享内存、异构加速器、集群多机部署与无服务器云函数等多元并行/分布式方案中进行严谨的架构选型与权衡论证。
    \item \textbf{[开发 (领域核心)]} 在前述八大算法领域中挑选其一，独立设计、实现、分析并评估一套针对特定物理与计算约束的高性能并发系统组件（例如：开发一个用于自动化仓储调度排程的并行遗传算法引擎）。
    \item \textbf{[评估 (领域核心)]} 深入审查并客观评判一个现存复杂并行组件（或由同学开发的并发系统），精准指出其在架构设计、同步机制、数据局部性或异常容错方面的深层次缺陷。
    \item \textbf{[评估 (领域核心)]} 针对同一计算基准算法（如大矩阵乘法），分别开发多核 CPU、分布式集群（MPI）以及 GPU（CUDA/OpenCL）三个不同版本的工程实现，严谨实测并全面对比三者在纯算力吞吐量与能耗效率维度的量化表现。
\end{enumerate}

\section{专业素养与思维倾向}

在并行与分布式计算领域的学习与工程实践中，重点培育学生的以下专业心智模型与素养特质：
\begin{itemize}
    \item \textbf{严密分析性与并发思维 (Analytical \& Concurrency Mindset)}：能够彻底跳出初学者习惯的单线程、单步线性确定性执行思维定式，建立交错时空、部分有序与状态空间爆炸下的非确定性并发系统全局观。
    \item \textbf{细致防御性与对部分失效的敬畏 (Meticulous \& Fault-tolerant Awareness)}：深刻理解并发竞态隐蔽性强、偶发性高以及分布式系统网络延迟不可预测、局部节点故障不可避免的现实规律，在接口设计与代码实现中始终秉持防卫性设计哲学，不心存侥幸。
    \item \textbf{权衡与工程折中意识 (Pragmatic \& Trade-off Oriented)}：深刻洞察分布式世界中的核心理论边界（如布鲁尔 CAP 定理、时延-一致性 PACELC 定理、阿姆达尔定律），能够在强一致性与高可用性、硬件成本与能耗极限、开发复杂度与运行效率之间做出清醒、理性的工程折中。
    \item \textbf{社会责任感与服务韧性考量 (Responsible \& Resilient)}：高度关注支撑现代社会运转的金融、交通、医疗与公共基础设施分布式系统的鲁棒性，时刻警惕单点故障与雪崩效应可能对公众生活造成的灾难性影响，坚守容错冗余、优雅降级与数据容灾的职业道德底线。
\end{itemize}

\section{数学基础要求}

并行与分布式计算依赖扎实的形式化数学工具与应用数学素养：
\begin{itemize}
    \item \textbf{离散数学与图论 (Discrete Mathematics \& Graph Theory)}：偏序集、哈斯图、拓扑排序；有向无环图（DAG，用于描述计算任务因果依赖图与工作流）；状态转移系统与自动机可达性分析。
    \item \textbf{数理逻辑与时态逻辑 (Logic \& Temporal Logic)}：命题逻辑与谓词演算；线性时态逻辑（LTL）与计算树逻辑（CTL），用于形式化表述与证明并发系统的安全性不变量（“坏事绝不发生”）与活性属性（“好事终将发生”）。
    \item \textbf{概率论、随机过程与排队论 (Probability \& Queuing Theory)}：离散与连续概率分布、随机变量期望与方差；泊松过程、马尔可夫链状态转移；排队论模型（M/M/1、Little 定律），用于定量分析通信信道排队延迟、吞吐瓶颈与缓存命中率。
    \item \textbf{线性代数与微积分 (Linear Algebra \& Calculus)}：向量与矩阵乘法、张量运算、特征值分解；多元微积分与偏微分方程离散化求解，支撑高性能数值计算与深度学习算子并行加速。
\end{itemize}

\section{交叉主题与跨学科共享概念}

\subsection{先修概念与依赖关系}
学习并行与分布式计算领域计算机核心知识的前提是熟练掌握以下先修概念：
\begin{itemize}
    \item \textbf{SDF (软件开发基础)}：程序静态代码与动态执行实体的区别；规约与实现的解耦；变量生命周期与内存表示；顺序控制流（条件分支、迭代循环）；过程/函数/方法调用机制与调用栈；数组与基础连续存储。
    \item \textbf{MSF (数学与统计学基础)}：命题逻辑；离散数学结构（有向图、状态转移机）；初等线性代数与微积分基础直觉。
    \item \textbf{FPL (编程语言理论基础)}：事件驱动架构；面向对象封装与多态；函数式编程范式、高阶函数与不可变数据模型。
    \item \textbf{SF (系统基础)}：分层抽象系统架构；冯·诺依曼计算机体系结构及其存储层次。
\end{itemize}

\subsection{知识领域交叉与重叠}
并行与分布式计算与计算机科学其他知识领域存在深度交织，部分拓展内容可在相关专业课中协同讲授：
\begin{itemize}
    \item \textbf{SF (系统基础)}：远程过程调用（RPC）底层机制、系统基准性能测试与多级存储层次评估。
    \item \textbf{NC (网络与通信)}：网络传输协议栈、Socket 编程接口、流量拥塞控制与拓扑路由。
    \item \textbf{OS (操作系统)}：内核级进程与线程调度、虚拟内存映射、系统调用陷阱、互斥锁与信号量实现、内核级故障隔离。
    \item \textbf{FPL (编程语言理论基础)}：并发内存模型形式化语义（如 Java/C++ 内存模型）、语言级并发原语、类型系统对数据竞态的静态杜绝。
    \item \textbf{SE (软件工程)}：高并发系统非功能性需求分析、微服务设计模式、并发故障回归测试。
    \item \textbf{SEC (网络空间安全)}：跨租户容器侧信道隔离、拜占庭容错、安全多方计算与区块链加密共识。
\end{itemize}

\section{课程体系规划与开课建议}

由于并行与分布式计算的计算机核心仅规定了 10 学时的通识要求，高校通常\textbf{不需要也不建议}单独开设一门仅覆盖 计算机核心 的极短必修课。联合工作组推荐以下三种典型开课与融合方案：

\subsection{方案一：全课程矩阵分散渗透方案}
将 10 学时的 计算机核心 知识有机拆解并注入至各核心骨干课程之中：
\begin{itemize}
    \item 注入《程序设计语言》：讲解函数式不可变性、并发闭包与轻量级线程启动；
    \item 注入《操作系统》：深化讲解多核进程/线程调度、互斥同步原语与死锁防范；
    \item 注入《计算机体系结构》：讲授多核缓存一致性协议（MESI）、伪共享与 GPU SIMD 向量化；
    \item 注入《计算机网络》与《数据库系统》：讲授 Socket 通信、RPC、事务 ACID、多版本并发控制（MVCC）与 Raft 共识。
\end{itemize}

\subsection{方案二：独立专业核心课程方案}
开设一门 48--64 学时的独立专业课程《并行与分布式计算系统》，完整统揽 计算机核心（10 学时）与 领域核心（26+ 学时），将多核编程、GPU 异构加速与分布式系统三大支柱融会贯通，配合高强度的工程实验（如基于 OpenMP/MPI 的并行矩阵加速、基于 Raft 的分布式键值数据库实现）。

\subsection{方案三：多核与高性能计算聚焦型课程范例}
针对偏向系统结构与高性能计算方向的高校，工作组提供了一套经典的多核导向开课样板：
\begin{enumerate}
    \item \textbf{程序与执行}：计算机核心 + 多核环境下的线程池、细粒度任务、指令级并行（ILP）。
    \item \textbf{通信机制}：计算机核心 + 现代多核微架构、多级 Cache 层次、伪共享规避、并发无锁数据结构。
    \item \textbf{协调与同步}：计算机核心 + 阻塞与非阻塞同步、投机执行、分治数据并行。
    \item \textbf{软件工程}：计算机核心 + 基于 DAG 的 Work-Span 跨度分析、性能调优剖析器实战。
    \item \textbf{算法与应用}：计算机核心 + 基于工业项目的并发流式数据处理与工作窃取任务调度器研发。
\end{enumerate}

\section{真实能力任务规范}

\begin{taskbox}[任务 1：基于已知算法实现高性能并行/分布式软件组件]
\textbf{【任务情境】}：针对已有明确数学或串行逻辑定义的核心算法（如大矩阵相乘、分治快速排序、PageRank 图计算或分布式哈希聚合），将其重构为能够在现代多核多节点环境中高效并发运转的高性能工程组件。\\
\textbf{【能力陈述】}：深刻理解跨越不同编程语言与底层计算框架（如 C++ 标准并发库、OpenMP、MPI、Java JUC 或 CUDA）的并行与分布式核心概念，并能准确、高质量地将其映射为鲁棒的代码实现。\\
\textbf{【能力领域】}：软件工程 \quad \textbf{【能力单元】}：系统开发\\
\textbf{【支撑知识单元】}：并行分布式/并行算法与应用、编程语言/并发与分布式计算\\
\textbf{【核心要求】}：计算机核心 \quad \textbf{【达标等级】}：\textbf{应用与开发}
\end{taskbox}

\begin{taskbox}[任务 2：通过引入并行或分布式架构突破串行系统性能瓶颈]
\textbf{【任务情境】}：某企业级核心业务系统面临数据量激增与高突发访问，原有单机串行组件成为全系统的严重性能瓶颈，亟需通过并行化或分布式横向扩展进行架构升级。\\
\textbf{【能力陈述】}：准确评估在何种场景下、通过何种并行或分布式技术能够有效提升吞吐量与降低延迟（并清醒识别不可并行化的阿姆达尔瓶颈）；精准捕捉并发化突破口，完成系统重构，并通过量化基准测试实测验证加速比与扩展收益。\\
\textbf{【能力领域】}：软件 / 系统 / 应用\\
\textbf{【能力单元】}：系统设计 / 开发 / 评估 / 性能优化\\
\textbf{【支撑知识单元】}：并行分布式/并发软件工程、编程语言/并发与分布式计算\\
\textbf{【核心要求】}：计算机核心 \quad \textbf{【达标等级】}：\textbf{评估与开发}
\end{taskbox}

\begin{taskbox}[任务 3：在保障系统关键语义与正确性前提下修订规约以支持分布式扩展]
\textbf{【任务情境】}：将传统基于单机强一致性假设的旧版软件规范拓展为支持网络分布式部署与异步并发处理时，必须全面修订软件架构规约与交互协议。\\
\textbf{【能力陈述】}：在松弛串行顺序约束或引入跨网络异步通信机制时，能够运用时态逻辑与不变量推理，严密论证并确保系统不会引入数据竞态、死锁、活锁或一致性破裂等破坏性隐患，确保原有系统核心功能与业务安全语义始终成立。\\
\textbf{【能力领域】}：软件 / 系统 / 应用\\
\textbf{【能力单元】}：架构设计 / 系统维护\\
\textbf{【支撑知识单元】}：并行分布式/并发软件工程、编程语言/并发与分布式计算\\
\textbf{【核心要求】}：计算机核心 \quad \textbf{【达标等级】}：\textbf{解释与评估}
\end{taskbox}

\begin{taskbox}[任务 4：针对复杂工程需求在多种并行与分布式架构方案间进行评估与选型]
\textbf{【任务情境】}：面对一个融合了海量批处理、实时低延迟在线查询与海量并发写入的复杂工业系统建设需求，需要在异构计算与分布式架构之间做出总体技术选型。\\
\textbf{【能力陈述】}：全面权衡多核共享内存、异构数据加速（GPU/SIMD）、反应式事件驱动架构、集群分布式计算与弹性云服务等不同范式的相对优劣、硬件成本、研发复杂度与吞吐极限，为具体的工程问题制定最优的系统架构选型方案。\\
\textbf{【能力领域】}：软件 / 系统 / 应用\\
\textbf{【能力单元】}：架构设计与方案评估\\
\textbf{【支撑知识单元】}：并行分布式/并行算法与应用、编程语言/并发与分布式计算\\
\textbf{【核心要求】}：计算机核心 \quad \textbf{【达标等级】}：\textbf{评估}
\end{taskbox}

\section{附录：核心主题与技能要求等级对照表}

\begin{longtable}{p{0.18\textwidth} p{0.48\textwidth} c c c}
\toprule
\textbf{知识单元} & \textbf{核心知识点与具体主题} & \textbf{认知等级} & \textbf{分类} & \textbf{学时} \\
\midrule
\endfirsthead
\toprule
\textbf{知识单元} & \textbf{核心知识点与具体主题} & \textbf{认知等级} & \textbf{分类} & \textbf{学时} \\
\midrule
\endhead
\bottomrule
\endfoot
\multirow{3}{*}{\makecell[l]{并行分布式/程序与执行}} 
& 并行动作定义、Happens-before 先行发生关系、数据独立性、一致性、物理执行位置与通信故障 & 解释 & 计算机核心 & 1 \\
& 并行动作的启动机制：编排放置、过程式派生、事件驱动反应式、依赖驱动启动 & 解释 & 计算机核心 & 1 \\
& 底层映射与物理执行机制：CPU 数据/指令级并行、GPU/SIMD 异构加速、多核调度、集群/云弹性编排与新兴量子/分子计算 & 解释 & 领域核心 & 2 \\
\midrule
\multirow{3}{*}{\makecell[l]{并行分布式/通信机制}} 
& 通信基础要素：介质（通道、共享内存、异构显存、数据存储）、内存一致性与缓存局部性、协同数据存储抽象 & 解释 & 计算机核心 & 2 \\
& 基于通信的并发编程：网络通道、套接字（Socket）、远程过程调用（RPC）及特定语言共享内存原语运用 & 开发 & 计算机核心 & 1 \\
& 高级通信属性与扩展：单播/组播/路由拓扑、排队论建模、高级通道流控/复用、顺序与获取-释放内存模型、DSM/RDMA 仿真、Raft/Paxos 容错共识 & 解释 & 领域核心 & 6 \\
\midrule
\multirow{3}{*}{\makecell[l]{并行分布式/协调与同步}} 
& 协调与同步基础：依赖动作控制（屏障、汇合、生产者-消费者）、死锁条件与推进活性、硬件原子指令与互斥锁 & 解释 & 计算机核心 & 2 \\
& 协调与同步编程实践：并发任务生命周期终止控制、锁与同步器防竞态实战、框架事务性 API 编程 & 开发 & 计算机核心 & 1 \\
& 高级协调扩展：无锁与无等待算法、读写锁/顺序锁变体、逻辑/向量时钟、投机执行与协同取消、信号量与条件变量、工作窃取任务调度 & 解释 & 领域核心 & 6 \\
\midrule
\multirow{3}{*}{\makecell[l]{并行分布式/软件工程}} 
& 软件需求规范：基于时态逻辑表达安全性与活性、吞吐量/时延/可用性/能耗与 SLA 度量指标 & 解释 & 计算机核心 & 1 \\
& 违规识别与修复：数据竞态错误、检查-执行原子性破裂、活锁与死锁排查修复 & 评估 & 计算机核心 & 1 \\
& 规范化与经验评估：线性一致性形式化规约、TLA+ 建模、DAG 工作量-跨度效率分析、ThreadSanitizer 竞态检测与压力测试 & 评估 & 领域核心 & 3 \\
\midrule
\multirow{3}{*}{\makecell[l]{并行分布式/算法与应用}} 
& 并行与分布式算法表达与实现：在主流语言与并发框架中落地线程派生、共享内存同步、套接字与 RPC & 开发 & 计算机核心 & 1 \\
& 核心计算范式与系统架构总览：多核共享内存、反应式、数据并行、集群、云计算、开放分布式系统的全景权衡 & 解释 & 计算机核心 & 1 \\
& 八大核心算法应用领域：并行线性代数、大规模数据处理、图算法与组合搜索、偏微分仿真、并行 SAT 求解、图形学渲染、工作窃取资源调度、高并发金融微服务 & 开发 & 领域核心 & 9 \\
\end{longtable}
